\section{Motor de Cálculo: Procesamiento de Señales y Ajuste de Curvas}
\label{subsec:motor_calculo}

Este componente constituye el núcleo matemático del software, y es responsable de obtener los parámetros normativos a partir de los arreglos temporales de tensiones discretas adquiridas por el \emph{hardware}.

El motor de cálculo ha sido diseñado exclusivamente para cumplir con las normas IEC 60060-1 \cite{iec60060_1} e IEC 61083-2-1 \cite{iec61083_2}, aplicando criterios de modularidad, que permiten su validación automatizada e independiente del resto de la aplicación.

\subsection{Flujo algorítmico del DSP}
\label{subsubsec:flujo_algoritmico}

El algoritmo del módulo matemático se estructura en función lo establecido en la norma del ensayo de impulso \cite{iec60060_1}.
Este proceso comprende desde el acondicionamiento inicial del vector de datos brutos hasta la obtención final de los parámetros característicos de la onda.

Las etapas se ejecutan secuencialmente de la siguiente manera:
\begin{enumerate}
    \item \textbf{Eliminación de Offset}: Se calcula el valor de tensión continua ($V_{\text{offset}}$) promediando estadísticamente una ventana de datos adquiridos ($N_{\text{pretrigger}}$) en el instante previo al disparo (``pre-trigger''), donde la curva registrada ($U_\text{e}[n]$) es nula.
    \begin{equation}
        \label{eq:calculo_offset}
        V_{\text{offset}} = \frac{1}{N_{\text{pretrigger}}} \sum_{n=1}^{N_{\text{pretrigger}}} U_\text{e}[n]
    \end{equation}
    El vector de tensión corregida se calcula como:
    \begin{equation}
        \label{eq:eliminacion_offset}
        U_{\text{corregida}}[n] = y[n] - V_{\text{offset}}
    \end{equation}
    \item \textbf{Normalización de Polaridad}: Dado que los impulsos pueden ser de polaridad positiva o negativa, el algoritmo detecta el sentido del impulso evaluando el signo de la muestra extrema:
    \begin{equation}
        \label{eq:buscar_polaridad}
        \text{polaridad} = \text{sign}(\text{arg max} |U_{\text{corregida}}[n]|)
    \end{equation}
    El vector de tensión corregida se multiplica por este término para simplificar el resto del procesamiento, asegurando que el análisis sea exclusivamente en el semiplano positivo:
    \begin{equation}
        \label{eq:normalizacion_polaridad}
        U_{\text{abs}}[n] = \text{polaridad} \cdot U_{\text{corregida}}[n]
    \end{equation}
    \item \textbf{Normalización Unitaria}: Se normaliza la amplitud al dominio unitario respecto a su valor de cresta.
    \begin{equation}
        \label{eq:onda_normalizada}
        U_{\text{norm}}[n] = \frac{U_{\text{abs}}[n]}{\max(U_{\text{abs}}[n])}
    \end{equation}
===============================================================================
    \item \textbf{Cálculo del desfasaje temporal}: El corrimiento temporal, $t_L$, entre el frente de onda del último impulso con el de referencia, se calcula como el promedio de las diferencias de los instantes de $t_{30}$, $t_{50}$ y $t_{80}$, correspondientes al \qty{30}{\percent}, \qty{50}{\percent} y \qty{80}{\percent} de $U_\text{t}$ respectivamente, de ambos impulsos.
    \begin{equation}
        t_L = \frac{\left( t_{30, \text{ref}} - t_{30, \text{i}} \right) + \left( t_{50, \text{ref}} - t_{50, \text{i}} \right) + \left( t_{80, \text{ref}} - t_{80, \text{i}} \right)}{3} 
    \end{equation}
    \item \textbf{Alineación del eje temporal}: Se obtiene el vector temporal alineado del último impulso, $t_{aligned}[n]$, partiendo de su vector temporal original $t[n]$, al que se lo desplaza la cantidad $t_L$.
    \begin{equation}
        \label{eq:eje_alineado}
        t_{aligned}[n] = t[n] + t_L
    \end{equation}
    \item \textbf{Clasificación del impulso}: Se calcula la diferencia absoluta punto a punto entre el último impulso y el de referencia, usando sus vectores normalizados en amplitud.
    \begin{equation}
        \label{eq:umbral_desviacion}
        | U_{\text{i,norm}}[n] - U_{\text{ref, norm}}[n] | > umbral
    \end{equation}
    Si supera el umbral, la onda es cortada y se almacena el punto de separación $idx_{deviation}$. Caso contrario, es plena y continúa el análisis principal.
        \item \textbf{Cálculo del factor de escala}: Se calcula como el promedio del intervalo comprendido entre el \qty{30}{\percent} y \qty{80}{\percent} de $U_{\text{corregida}}[n]$ del frente de la última onda y la de referencia.
    \begin{equation}
        \overline{U_\text{i}} = \frac{1}{N} \sum_{n = k_{30\%}}^{k_{80\%}} U_\text{i}[n]
    \end{equation}
    \begin{equation}
        \overline{U_\text{ref}} = \frac{1}{N} \sum_{n = k_{30\%}}^{k_{80\%}} U_\text{ref}[n]
    \end{equation}
    \begin{equation}
        \label{eq:factor_escala_e}
        E = \frac{\overline{U_\text{i}}}{\overline{U_\text{ref}}}
    \end{equation}
    \item \textbf{Escalado de la Curva Base}: Debido a que la descarga disruptiva hace colapsar el ajuste por el corte abrupto de la señal, se sintetiza la curva base de la onda de referencia $U_\text{b, ref}[n]$, y se escala en amplitud con el factor de escala $E$.
    \begin{equation}
        \label{eq:curva_base_lic}
        U_\text{b, i}[n] = E \cdot U_\text{b, ref}[n]
    \end{equation}
    \item \textbf{Tiempo de corte}: Se analiza la cola del impulso, identificando el punto de máxima pendiente de caída mediante la derivada primera
    \begin{equation}
        \label{eq:gradiente_colapso}
        n_\text{abrupto} = \text{arg min}\left(\frac{\text{d} U_{\text{cola}}[n]}{\text{d} t}\right)
    \end{equation}
    y el punto de inflexión mediante la derivada segunda
    \begin{equation}
        \label{eq:derivada2_colapso}
        n_\text{inflexion} = \text{arg min}\left(\frac{\text{d}^2 U_\text{cola}[n]}{\text{d}^2 t}\right)
    \end{equation}
    para obtener la tensión del instante del colapso
    \begin{equation}
        U_\text{colapso} = U_\text{corregida}[n_\text{pico} + n_\text{inflexion}]
    \end{equation}
    Se calculan los instantes correspondientes al \qty{70}{\percent} y al \qty{10}{\percent} de $U_\text{colapso}$ en la descarga disruptiva, para extrapolar el instante de corte:
    \begin{equation}
        t_c = t_{70} + (U_{\text{colapso}} - U_{70}) \left(\frac{t_{10} - t_{70}}{U_{10} - U_{70}}\right)
    \end{equation}
    Finalmente, se calcula el tiempo de corte respecto al origen virtual:
    \begin{equation}
        T_c = t_c - O_1
    \end{equation}

===============================================================================

    \item \textbf{Segmentación de Datos}: Se selecciona una ventana de datos en el intervalo comprendido entre el \qty{20}{\percent} en el flanco de subida y el \qty{40}{\percent} en la cola de descarga del impulso.
    \item \textbf{Ajuste de curva base}: Se realiza un ajuste paramétrico no lineal sobre la ventana de muestras, utilizando el algoritmo de \textbf{Levenberg-Marquardt} \cite{iec60060_1}, teniendo como modelo de estimación la función de doble exponencial:
\begin{equation}
    \label{eq:modelo_doble_exponencial}
    u_d[n] = U \cdot \left( e^{-\frac{t - t_d}{\tau_1}} - e^{-\frac{t - t_d}{\tau_2}} \right)
\end{equation}
Donde los coeficientes desconocidos a determinar son:
\begin{itemize}
    \item $U$ el parámetro de amplitud escalar.
    \item $\tau_1$ la constante de tiempo de descarga.
    \item $\tau_2$ la constante de tiempo de carga.
    \item $t_d$ el retardo temporal del origen virtual.
\end{itemize}
    \item \textbf{Cálculo de Curva Residual}: Se obtiene como la diferencia algebraica entre la curva corregida ($U_{\text{corregida}}[n]$) y la curva base ($U_\text{b}[n]$) construída con los parámetros obtenidos del ajuste:
    \begin{equation}
        \label{eq:curva_residual}
        R[n] = U_{\text{corregida}}[n] - U_\text{b}[n]
    \end{equation}
    \item \textbf{Filtrado Digital de Fase Cero}: 
    Se obtiene la curva residual filtrada ($R_f[n]$), aplicando un filtro digital IIR a la curva residual ($R[n]$). Para reducir el ruido de muy alta frecuencia relacionado a la interferencia electromagnética, sin perder oscilaciones propias del ensayo.
    \item \textbf{Curva de Tensión de Ensayo}: Se sintetiza la curva de ensayo ($U_\text{t}[n]$) sumando algebraicamente la curva base analítica con la residual filtrada:
    \begin{equation}
        \label{eq:curva_ensayo}
        U_\text{t}[n] = U_\text{b}[n] + R_f[n]
    \end{equation}
    \item \textbf{Cálculo de parámetros característicos}: A partir de la curva de tensión de ensayo $U_\text{t}[n]$, se calculan los parámetros normativos \cite{iec60060_1}. Dado que la señal procesada es de tiempo discreto, hay valores analíticos que pueden no corresponder con las muestras obtenidas, por lo que para encontrarlos se realiza una interpolación lineal.
    \begin{itemize}
        \item \textbf{Valor de cresta} ($U_\text{t}$): Es la amplitud máxima de la curva de tensión de ensayo $U_\text{t}[n]$:
        \begin{equation}
            \label{eq:valor_cresta}
            U_\text{t} = \text{max}(U_\text{t}[n])
        \end{equation}
        \item \textbf{Tiempo de Frente} ($T_1$): Se calcula como el cociente del intervalo temporal transcurrido entre el instante $t_{30}$ (amplitud correspondiente al \qty{30}{\percent} de $U_\text{t}$) y el instante $t_{90}$ (amplitud correspondiente al \qty{90}{\percent} de $U_\text{t}$), dividiendo por $\num{0.6}$:
        \begin{equation}
            \label{eq:tiempo_frente}
            T_1 = \frac{t_{90} - t_{30}}{\num{0.6}}
        \end{equation}
        \item \textbf{Origen Virtual} ($O_1$): El inicio ficticio de la descarga eléctrica se calcula:
        \begin{equation}
            \label{eq:origen_virtual}
            O_1 = t_{30} - \num{0.5} \cdot T_{ab}
        \end{equation}
        \item \textbf{Tiempo de Cola} ($T_2$): El intervalo temporal entre el origen virtual $O_1$ y el instante $t_{50}$ (amplitud correspondiente al \qty{50}{\percent} de $U_\text{t}$) en el flanco descendente de la onda, se calcula:
        \begin{equation}
            \label{eq:tiempo_cola}
            T_2 = t_{50} - O_1
        \end{equation}
        \item \textbf{Sobreimpulso Relativo} ($\text{OS}\%$): Se calcula como la diferencia entre el valor extremo de la curva registrada ($U_\text{e}$):
        \begin{equation}
            \label{eq:max_curva_registrada}
            U_\text{e} = \max(U_\text{e}[n])
        \end{equation}
        y el valor máximo de la curva base ($U_\text{b}$):
        \begin{equation}
            \label{eq:max_curva_base}
            U_\text{b} = \max(U_\text{b}[n])
        \end{equation}
        expresado en porcentaje:
        \begin{equation}
            \label{eq:sobreimpulso_relativo}
            \text{OS}\% = \frac{U_\text{e} - U_\text{b}}{U_\text{e}} \cdot \num{100}
        \end{equation}
    \end{itemize}
\end{enumerate}